% Standalone source-record documentation; no research attempts.
\documentclass[11pt]{article}
\usepackage[margin=25mm]{geometry}
\usepackage{fontspec,unicode-math}
\setmainfont{Latin Modern Roman}
\setmathfont{Latin Modern Math}
\usepackage{amsmath,amssymb,mathtools,xurl,hyperref}
\hypersetup{colorlinks=true,urlcolor=blue,linkcolor=blue}
\setcounter{secnumdepth}{0}
\setlength{\parindent}{0pt}
\setlength{\parskip}{5pt}
\setlength{\emergencystretch}{3em}
\newcommand{\sref}[2]{\hyperlink{source:#1:#2}{[S#2]}}
\newcommand{\cataloguescope}[1]{\paragraph{Source record.}#1}
\newcommand{\floor}[1]{\left\lfloor#1\right\rfloor}
\begin{document}
% UnsolvedMath ID 30005937; source catalogue code OWR-14298374-006
\setcounter{section}{30005936}
\section{Autoencoder Optimization Perspectives}\label{30005937}
{\small\sffamily UnsolvedMath ID 30005937 · Computer Science}
\subsection{Definitions and mathematical statement}

\paragraph{Shared notation.}Standard number systems are $\mathbb N,\mathbb Z,\mathbb Q,\mathbb R,\mathbb C$. Specialized notation and index ranges follow the cited source and the definitions below.


An autoencoder consists of an encoder and decoder, fitting reconstruction $x\mapsto f_{\mathrm{dec}}(f_{\mathrm{enc}}(x))$. Reconstruction loss compares the decoded result with the original input. Conditional expectation is an optimal squared-error predictor for a specified sigma-algebra.
\paragraph{Record scope and qualifications.}
The catalogue title is retained as metadata. Named certificates, contracts, bounds or frameworks in that title are not asserted as results of the original source extract.
The supplied registry flags this extraction as unrecoverable. It is displayed as a literal source fragment; the definitions provide context rather than a replacement formal statement.

This is incomplete expository text, without a distinct fully specified unresolved proposition.

Source-record limitation: Source flags the statement as unrecoverable: The official report describes three equivalent optimization viewpoints and their modeling consequences; it does not state an unresolved mathematical problem.
The source text below is rendered literally where necessary; damaged formulas and composite extracts are not silently repaired into a new mathematical claim.
The supplied catalogue marks this record as unrecoverable. Its extract and terminology are documented, but a complete original problem has not been reconstructed.
\paragraph{Statement recorded in the supplied source.}
In particular, there is some duality in the way the minimization over f \textbackslash{}in F is performed, as one can decide to either (i) simultaneously minimize over fenc and fdec , which is the standard way to proceed when training neural networks; (ii) minimize first over the encoder part, which allows to reformulate the min- imization as the well-known problem of finding principal manifolds; (iii) minimize first over the decoder part, which is natural when thinking of the reconstruction error as some total variance to be decomposed using a conditioning on the values of the encoder. 1.2. Reformulating autoencoders with conditional expectations. We dis- cuss here how to reformulate the training of autoencoders with conditional expec- tations, and provide alternative interpretations to the reconstruction error. We minimize the reconstruction error (3) by first minimizing over the decoder part for a given encoding function, as already considered in [7]. This approach is natural in molecular dynamics, as it is reminiscent of free energy computations [3, 11] where average quantities are computed for a fixed value of the collective variable fenc . From a mathematical viewpoint, it corresponds to introducing conditional averages associated with fixed values of the encoder. The loss function for unsupervised least-squares can be rewritten as h i h i 2 2 inf E kX − f (X)k = inf inf E kX − fdec ◦ fenc (X)k f \textbackslash{}inF fenc \textbackslash{}inFenc fdec \textbackslash{}inFdec h i 2 (4) = inf E X − gf⋆enc ◦ fenc (X) , fenc \textbackslash{}inFenc \sref{30005937}{2}
\paragraph{Source status.}
The supplied catalogue labels this source record open; that metadata was checked against the supplied snapshot on 2026-10-11. The label is preserved for provenance and is not an independent verification of present open status or mathematical correctness. The catalogue literature assessment is dated 2026-08-22.
\subsection{Short English statement}
The fragment describes several ways to optimize or reformulate an autoencoder.
\subsection{Sources}
\begin{itemize}
\item[\hypertarget{source:30005937:1}{S1}] Ulam AI and the UnsolvedMath Contributors, UnsolvedMath: A Curated Collection of Open Mathematics Problems, record 30005937 (OWR-14298374-006). Catalogue curation and metadata: CC BY 4.0. The source extract is retained; definitions, English summaries and qualifications are editorial additions.
\url{https://huggingface.co/datasets/ulamai/UnsolvedMath}.
\item[\hypertarget{source:30005937:2}{S2}] Constrained Dynamics, Stochastic Numerical Methods and the Modeling of Complex Systems (2024), Oberwolfach Reports, DOI 10.4171/owr/2024/26, pp. 1445--1516
\url{https://doi.org/10.4171/owr/2024/26}.
\item Source-specific Oberwolfach Reports licences, including CC BY-SA 4.0 where applicable, are retained; see THIRD\_PARTY\_NOTICES.md.
\end{itemize}
\paragraph{Preparation and review.}
The original source and UnsolvedMath contributors retain attribution for the source extract. Codex (GPT-6; exact serving revision not exposed) assisted with definitions, English summaries, scope notes and formatting. No new human mathematical authorship or independently verified solution is asserted. Permanent source ID 30005937 and code OWR-14298374-006 checked against the complete supplied registry and current repository. Distinct upstream records, including repeated questions and broad or damaged extracts, are retained for complete coverage of the source CS category. A repeated record is not a new independent conjecture. Complete mathematical-equivalence and current-literature review has not been performed. Preparation and final review follow contribution prompts C and D, with source gaps disclosed. Structural and rendering checks do not establish a new theorem or open status.
\label{end:30005937}
\end{document}
